\section{几何类域论}\srcpage{4}

设 $X$ 是连通、射影、光滑的代数曲线（紧连通 Riemann 面），亏格为 $g$。取定基点 $x_0$、自由 $\Z$-模
\[
  H_1(X,\Z)=\pi_1^{\mathrm{ab}}(X):=\pi_1(X,x_0)/[\pi_1(X,x_0),\pi_1(X,x_0)]
\]
的一组基 $(\delta_1,\dots,\delta_{2g})$，以及复向量空间 $H^{0}(X,\OmX)$ 的一组基 $(\omega_1,\dots,\omega_g)$。周期
\[
  \Pi_i=\left(\int_{\delta_i}\omega_1,\dots,\int_{\delta_i}\omega_g\right)\in\C^{g},
  \qquad i=1,\dots,2g,
\]
\srcpage{5}在 $\Z$ 上线性无关，并生成一个格 $\Lambda\subset\C^{g}$。$X$ 的雅可比簇是 $g$ 维复环面
\[
  J(X)=\C^{g}/\Lambda.
\]
有一个解析映射
\[
  \varphi:X\to J(X),\quad
  x\mapsto\left(\int_{x_0}^{x}\omega_1,\dots,\int_{x_0}^{x}\omega_g\right)+\Lambda,
\]
它把 $x_0$ 送到 $J(X)$ 的单位元 $0$。

设 $\Pic(X)$ 是参数化 $X$ 上线丛同构类的 Picard 概形。线丛的张量积给出群结构，使它成为群概形，并落在正合列
\[
  1\to\Pic^{0}(X)\to\Pic(X)\xrightarrow{\deg}\Z\to 0
\]
中；这里 $\deg(\cL)$ 是线丛 $\cL$ 的次数，$\Pic(X)$ 的单位分支 $\Pic^{0}(X)$ 是阿贝尔簇。

\begin{theorem}[Abel--Jacobi]
雅可比簇 $J(X)$ 是阿贝尔簇 $\Pic^{0}(X)$ 的底层复环面，解析态射
\[
  \varphi:X\to J(X)\cong\Pic^{0}(X)
\]
正是把 $x\in X$ 送到 $\cO_X([x]-[x_0])$ 在 $\Pic^{0}(X)$ 中的同构类的代数态射。
\end{theorem}

由 $\varphi$ 诱导的同态
\[
  \pi_1(X,x_0)\to\pi_1(J(X),0)
\]
是商映射
\[
  \pi_1(X,x_0)\twoheadrightarrow
  \pi_1(X,x_0)/[\pi_1(X,x_0),\pi_1(X,x_0)]=\Lambda.
\]
所以对每个特征标 $\chi:\pi_1(X,x_0)\to\C^{\times}$，存在唯一的特征标
$\Aut_{\chi}:\pi_1(\Pic^{0}(X),0)\to\C^{\times}$，使 $\chi=\Aut_{\chi}\circ\varphi_{*}$。

流形上的局部系统（$\C$-向量空间）无非是该流形基本群的有限维 $\C$-表示。因此我们用解析方法证明了：

\begin{theorem}
对 $X$ 上任意秩 $1$ 的局部系统 $E$，存在 $\Pic^{0}(X)$ 上唯一的秩 $1$ 局部系统 $\Aut_E$（在原点刚性化），使 $\varphi^{*}\Aut_E=E$。此外，若 $m:\Pic^{0}(X)\times\Pic^{0}(X)\to\Pic^{0}(X)$ 是群运算，则有典范同构
\[
  m^{*}\Aut_E\cong\Aut_E\boxtimes\Aut_E.
\]
\end{theorem}

\srcpage{6}
\begin{proof}[定理 1.2 的代数证明]
对每个整数 $d\ge 0$，对称群 $\mathfrak{S}_d$ 以显然的方式作用在积
\[
  X^{d}=\overbrace{X\times\cdots\times X}^{d}
\]
上。

商概形
\[
  X^{(d)}=X^{d}/\mathfrak{S}_d
\]
连通、射影、光滑，维数为 $d$；它是 $X$ 上有效除子 $D=\sum_x d_x[x]$（$d_x\ge 0$）的模空间，这些除子的次数为 $\sum_x d_x=d$。态射 $\varphi$ 诱导态射
\[
  \varphi^{(d)}:X^{(d)}\to\Pic^{0}(X),\quad D\mapsto\cO_X(D-d[x_0]).
\]
只要 $d\ge 2g-1$，$\varphi^{(d)}$ 就是秩 $d-g$ 的射影丛，因此其纤维单连通。

对 $X$ 上每个秩 $n$ 的局部系统 $E$ 与每个整数 $d\ge 0$，外张量积 $E^{\boxtimes d}$ 是 $X^{d}$ 上秩 $n^{d}$ 的局部系统，带有 $\mathfrak{S}_d$ 的作用，该作用提升 $X^{d}$ 上的作用。记商态射为 $r:X^{d}\twoheadrightarrow X^{d}/\mathfrak{S}_d=X^{(d)}$，则 $\C$-向量空间的可构造层 $r_{*}E^{\boxtimes d}$ 带有 $\mathfrak{S}_d$ 的作用。记
\[
  E^{(d)}=(r_{*}E^{\boxtimes d})^{\mathfrak{S}_d}\subset E^{\boxtimes d}
\]
为该作用的不动向量子层。可以验证，$E^{(d)}$ 在 $D\in X^{(d)}$ 处的纤维等于
\[
  \tag*{(1.3)} (E^{(d)})_D=\bigotimes_x\Sym^{d_x}E_x.
\]

若现在 $E$ 是 $X$ 上秩 $1$ 的局部系统，则 $E^{(d)}$ 实际上是 $X^{(d)}$ 上秩 $1$ 的局部系统。对每个 $d\ge 2g-1$，$E^{(d)}$ 因此在态射 $\varphi^{(d)}:X^{(d)}\to\Pic^{0}(X)$ 的单连通纤维上常值，并下降为 $\Pic^{0}(X)$ 上的局部系统，正是 $\Aut_E$。
\end{proof}

这是处处非分歧的 Drinfeld--Langlands 对应中的 $\mathrm{GL}(1)$ 情形。
